% !PRINTED 245
\weberpage{245}
%<<CONT
aber nicht $\equiv 0\ (\mathrm{mod}\ p)$ sein; denn sonst liesse sich aus (6) eine Congruenz ableiten:

\[
P \equiv T P_{\nu-1}\ (\mathrm{mod}\ p),
\]

in der $T$ eine nicht constante Function von $t$ ist, und $P$ wäre nicht irreducibel nach dem Modul $p$.

Ist nun $Q$ irgend eine ganze Function von $t$, die der Bedingung

\[
Q P_1 \equiv 0\ (\mathrm{mod}\ P, p)
\]

genügt, so folgt aus (6) durch Multiplication mit $Q$:

\[
Q P_2 \equiv 0,\ Q P_3 \equiv 0,\ \ldots,\ Q P_\nu \equiv 0\ (\mathrm{mod}\ P, p),
\]

also $Q \equiv 0$. Und wenn wir also für $P_1$, $Q$ die Reste der Functionen $F(t)$, $F_1(t)$ setzen, so ist hiermit der Satz 1. bewiesen.

Die Reste aller Functionen $F(t)$ nach dem Modul $P$ sind von der Form

\weq{7}{X = X(t) = x_0 + x_1 t + \cdots + x_{n-1} t^{n-1},}

und wenn man nur die nach dem Modul $p$ incongruenten unter ihnen haben will, so genügt es, die Co\"efficienten $x_0, x_1, \ldots x_{n-1}$ die Reihe der Zahlen $0, 1, 2, \ldots, p-1$ durchlaufen zu lassen.

\gesperrt{Es giebt also $p^n$ und nicht mehr nach den Moduln $P$, $p$ incongruente Functionen.}

\S.\,64.

\gesperrt{Congruenzkörper.}

Es ist nur eine andere Ausdrucksweise für die im vorigen Paragraphen durchgeführten Betrachtungen, wenn man mit \gesperrt{Galois} eine imaginäre Zahl $\varepsilon$ einführt, die der Congruenz

\weq{1}{P(\varepsilon) \equiv 0\ (\mathrm{mod}\ p)}

genügt. Eine ganze rationale Zahl dieser Art giebt es nicht, sobald $n > 1$ ist, weil sonst $P(t) - P(\varepsilon) = (t-\varepsilon)Q$ durch $t-\varepsilon$ theilbar, mithin $P(t) \equiv (t-\varepsilon)Q$, und $P(t)$ nicht irreducibel wäre. Gleichwohl kann man, wie man in der gewöhnlichen Zahlenlehre die imaginäre Einheit $i = \sqrt{-1}$ einführt, ein solches Symbol $\varepsilon$ in der Rechnung benutzen. Man rechnet damit, sofern es sich um Addition, Subtraction und Multiplication handelt, nach den Regeln der Buchstabenrechnung, und kann dabei nach Belieben die Gleichung $P(\varepsilon) = 0$ benutzen. Man kann sogar unter $\varepsilon$ geradezu eine Wurzel der Gleichung $P(x) = 0$
%CONT>>
